\section{Adaptation To Transactional Memory}
\subsection{Introducing Generalisation}
\label{sec:adaptation_to_transactional_memory}

\begin{BLACKSHEET}

In this model, locating a target node requires a top-down traversal where a thread must observe all active transactions within its path. Formally, ensuring that an update to a set is visible to all overlapping operations is non-trivial; while visibility is inherent for supersets of this set, it cannot be guaranteed for arbitrary subsets without a common synchronization point. Since no single element can consistently represent all minimal subsets (e.g., disjoint leaves), a global consensus on the frame's state is required. This formal constraint justifies reliance on localized frames and edge-based indirection as the primary mechanism for maintaining a consistent view across concurrent operations.

Consequently, we bypass the "lowest common ancestor" (LCA) approach for frame identification, as it fails to guarantee atomicity during complex rebalancing to the abstraction of frame sets. Instead, we encapsulate within the descriptor all memory cell data accessed by the current operation. This ensures that all modifications to nodes and edges are performed strictly inside this descriptor. The primary data structure positions will store exclusively references, as their values no longer remain independently valid and are not updated atomically while an operation is in progress over them.

This two-phase approach—logical commit followed by physical propagation—mirrors the Multi-word Compare-and-Swap (MCAS) semantics proposed by Harris et al. \cite{PracticalMCAS}, ensuring that updates to multiple memory locations remain linearizable through a single shared descriptor. By its very nature, this mechanism converges toward the software transactional memory model, demonstrating that the structural primitives required to guarantee lock-free properties in search trees can be readily transformed into a universal construction, with the only additional overhead being descriptor management. Furthermore, this design inherently preserves disjoint-access parallelism, ensuring that operations targeting non-overlapping memory regions execute independently. This property yields significant performance advantages over traditional lock-based STM implementations.

Based on this descriptor-centric architecture, we decompose each transaction into three atomic phases that ensure safe state transitions across the frame:
\begin{itemize}
    \item \textbf{Acquire.} The thread attempts to install its descriptor $D$ into each target memory cell $R_j \in \{r_1, \dots, r_n\}$ via $\text{CAS}(R_j, \text{old}_j, D)$. If a cell is already held by another descriptor $D'$, the thread must first invoke a recursive helping routine to assist the conflicting operation. This ensures that acquisition is never blocked by stalled processes. Furthermore, descriptors must be installed in an ascending order based on memory addresses to prevent deadlocks.
    \item \textbf{Decision.} The logical linearization point occurs here. Once all target cells are successfully acquired, the thread attempts an atomic CAS on the descriptor's status field, transitioning it from \texttt{Undecided} to either \texttt{Committed} or \texttt{Aborted}. The success of this single update determines the outcome of the entire multi-word operation.
    \item \textbf{Release.} In the final cleanup phase, the thread resets each node's state to a concrete value. If the status is \texttt{Committed}, the cell is updated to $\text{new}_j$; otherwise, it is reverted to $\text{old}_j$. This phase effectively releases the descriptor-based ownership and makes the structural change visible to the rest of the system.
\end{itemize}

On the one hand, this approach closely aligns with the implementation algorithm proposed in Efficient Multi-Word Compare and Swap \cite{EfficientMCAS}. The proof utilized in that work relies on theorems and lemmas to demonstrate the framework's correctness. In contrast, our objective is to present a formal specification of this algorithm, optimized specifically for reader threads that require a guarantee of data immutability until the operation completes. Writing and proving such specifications mathematically is highly convenient, and their correctness can be rigorously verified using the TLA+ language and the TLC model checker.

On the other hand, this approach follows the lineage of descriptor-based atomicity and was later refined in composable memory models \cite{ComposableMT}. However, our model fundamentally alters the semantics of read operations. To guarantee a lock-free view, any process accessing a cell must proactively resolve any pending descriptor. This dependency effectively elevates a simple load into a ``heavyweight'' operation, as a read may now trigger a physical write to memory. While this prevents the observation of transient states, it imposes a non-trivial overhead: every read becomes a potential state transition in the global descriptor space, complicating the performance profile in read-heavy workloads.

\end{BLACKSHEET}
